\documentclass[10pt,a4paper]{amsart}
\usepackage{amsmath,amssymb,amsthm,graphicx}
% Retained preamble for the Stacks Project chapter morphisms.tex (Morphisms of Schemes).
% The chapter file pinned-original-source.tex carries only \input{preamble} on its first line, so no
% preamble text lives inside the chapter itself.  The Stacks Project uses one single shared preamble for
% every chapter; the snapshot stages it verbatim in sources/stacks-chow/preamble.tex.  The block below is
% copied mechanically, line for line, from that file's lines 14-18 (the Xy-pic setup) and lines 154-217
% (the native theorem declarations together with the mathematical macros).  Nothing was added, removed or
% redefined.  The wrapper class and the amsmath/amssymb/amsthm/graphicx base packages are supplied by the
% assembler, so the class/\externaldocument part of the shared preamble is not repeated here.
% For commutative diagrams we use Xy-pic
\usepackage[all]{xy}

% We use 2cell for 2-commutative diagrams.
\xyoption{2cell}
% Theorem environments.
%
\theoremstyle{plain}
\newtheorem{theorem}[subsection]{Theorem}
\newtheorem{proposition}[subsection]{Proposition}
\newtheorem{lemma}[subsection]{Lemma}

\theoremstyle{definition}
\newtheorem{definition}[subsection]{Definition}
\newtheorem{example}[subsection]{Example}
\newtheorem{exercise}[subsection]{Exercise}
\newtheorem{situation}[subsection]{Situation}

\theoremstyle{remark}
\newtheorem{remark}[subsection]{Remark}
\newtheorem{remarks}[subsection]{Remarks}

\numberwithin{equation}{subsection}

% Macros
%
\def\lim{\mathop{\mathrm{lim}}\nolimits}
\def\colim{\mathop{\mathrm{colim}}\nolimits}
\def\Spec{\mathop{\mathrm{Spec}}}
\def\Hom{\mathop{\mathrm{Hom}}\nolimits}
\def\Ext{\mathop{\mathrm{Ext}}\nolimits}
\def\SheafHom{\mathop{\mathcal{H}\!\mathit{om}}\nolimits}
\def\SheafExt{\mathop{\mathcal{E}\!\mathit{xt}}\nolimits}
\def\Sch{\mathit{Sch}}
\def\Mor{\mathop{\mathrm{Mor}}\nolimits}
\def\Ob{\mathop{\mathrm{Ob}}\nolimits}
\def\Sh{\mathop{\mathit{Sh}}\nolimits}
\def\NL{\mathop{N\!L}\nolimits}
\def\CH{\mathop{\mathrm{CH}}\nolimits}
\def\proetale{{pro\text{-}\acute{e}tale}}
\def\etale{{\acute{e}tale}}
\def\QCoh{\mathit{QCoh}}
\def\Ker{\mathop{\mathrm{Ker}}}
\def\Im{\mathop{\mathrm{Im}}}
\def\Coker{\mathop{\mathrm{Coker}}}
\def\Coim{\mathop{\mathrm{Coim}}}

% Boxtimes
%
\DeclareMathSymbol{\boxtimes}{\mathbin}{AMSa}{"02}

%
% Macros for moduli stacks/spaces
%
\def\QCohstack{\mathcal{QC}\!\mathit{oh}}
\def\Cohstack{\mathcal{C}\!\mathit{oh}}
\def\Spacesstack{\mathcal{S}\!\mathit{paces}}
\def\Quotfunctor{\mathrm{Quot}}
\def\Hilbfunctor{\mathrm{Hilb}}
\def\Curvesstack{\mathcal{C}\!\mathit{urves}}
\def\Polarizedstack{\mathcal{P}\!\mathit{olarized}}
\def\Complexesstack{\mathcal{C}\!\mathit{omplexes}}
% \Pic is the operator that assigns to X its picard group, usage \Pic(X)
% \Picardstack_{X/B} denotes the Picard stack of X over B
% \Picardfunctor_{X/B} denotes the Picard functor of X over B
\def\Pic{\mathop{\mathrm{Pic}}\nolimits}
\def\Picardstack{\mathcal{P}\!\mathit{ic}}
\def\Picardfunctor{\mathrm{Pic}}
\def\Deformationcategory{\mathcal{D}\!\mathit{ef}}

% --- appended 2026-10-12 (Chief): the project own definitions of the reference/slogan/history
% display environments, copied mechanically from the shared Stacks preamble lines 8-12
% (\usepackage{verbatim} plus the three \newenvironment lines, which the chapter preamble keeps
% verbatim from the project preamble). Nothing added, removed or redefined: these environments
% carry no mathematical content and the project itself typesets them as comments. Without these
% three lines a counted statement whose own text uses them fails with Environment slogan undefined.
\usepackage{verbatim}
\newenvironment{reference}{\comment}{\endcomment}
\newenvironment{slogan}{\comment}{\endcomment}
\newenvironment{history}{\comment}{\endcomment}
\usepackage{cleveref}
\crefname{originallemma}{original lemma}{original lemmas}
\Crefname{originallemma}{Original lemma}{Original lemmas}
\makeatletter
\newlabel{morphisms-lemma-generically-finite}{{02NW}{source}{}{sourcepointer-1}{}}
\newlabel{morphisms-lemma-generically-finite@cref}{{[originallemma][0][]02NW}{[1][0][]source}{}{}{}}
\makeatother
\crefname{originallemma}{original lemma}{original lemmas}
\Crefname{originallemma}{Original lemma}{Original lemmas}
\makeatletter
\newlabel{algebra-lemma-generically-finite}{{02ML}{source}{}{sourcepointer-2}{}}
\newlabel{algebra-lemma-generically-finite@cref}{{[originallemma][0][]02ML}{[1][0][]source}{}{}{}}
\makeatother
\crefname{originallemma}{original lemma}{original lemmas}
\Crefname{originallemma}{Original lemma}{Original lemmas}
\makeatletter
\newlabel{lemma-finite-fibre}{{02NG}{source}{}{sourcepointer-3}{}}
\newlabel{lemma-finite-fibre@cref}{{[originallemma][0][]02NG}{[1][0][]source}{}{}{}}
\makeatother
\crefname{originallemma}{original lemma}{original lemmas}
\Crefname{originallemma}{Original lemma}{Original lemmas}
\makeatletter
\newlabel{properties-lemma-maximal-points-affine}{{01ZV}{source}{}{sourcepointer-4}{}}
\newlabel{properties-lemma-maximal-points-affine@cref}{{[originallemma][0][]01ZV}{[1][0][]source}{}{}{}}
\makeatother
\crefname{originallemma}{original lemma}{original lemmas}
\Crefname{originallemma}{Original lemma}{Original lemmas}
\makeatletter
\newlabel{lemma-quasi-compact-generic-point-not-in-image}{{02NE}{source}{}{sourcepointer-5}{}}
\newlabel{lemma-quasi-compact-generic-point-not-in-image@cref}{{[originallemma][0][]02NE}{[1][0][]source}{}{}{}}
\makeatother
\title{A quasi-finite morphism is generically finite}
\author{The Stacks Project Authors}
\date{stacks-project:morphisms.tex:local-snapshot-20261012; GNU FDL 1.2 or later}
\begin{document}
\maketitle
\setcounter{subsection}{1}
\noindent\textit{Editorial note.} All delivered mathematical text is the author's own, copied from the pinned Stacks Project chapter file morphisms.tex; no statement or proof step was rewritten, reordered or removed.\par\medskip
\noindent\textit{Editorial note.} The chapter file pinned-original-source.tex carries only \textbackslash{}input\{preamble\} on its first line, so no preamble text lives inside the chapter. The retained file sources/stacks-morphisms/retained-preamble.tex is a mechanical, line-for-line copy of the project's single shared preamble as staged in the snapshot. Nothing was added, removed or redefined, and the excerpt serves only as preamble, never as a span coordinate.\par\medskip
\noindent\textit{Editorial note.} Both counted spans start at the \textbackslash{}begin line of their environment, so the statement and the author's proof are delivered whole; surrounding prose of the section is deliberately outside the counted spans and is carried as standard prerequisites instead.\par\medskip
\noindent\textit{Editorial note.} The citation fields (source\_id, source\_version, authors, source\_url, license.name, source\_title) are identical to those of every other card this lane has produced from this chapter, so the batch selector treats them as one source.\par\medskip
\noindent\textit{Editorial note.} Read in full under the Lead's 12:4x rule (screening hit re-read, pb 2277B). The screening hit was the phrase 'nontrivial specialization', where 'trivial' is part of a mathematical term, not an omission; the proof is complete.\par\medskip
\noindent\textit{Allowed standard prerequisites (editorial).}\begin{itemize}
\item Local quasi-finiteness bounds the fibres above points. Lemma \textbackslash{}ref\{lemma-quasi-finite-fibre\} and the surrounding discussion, used at proof lines 13230-13250. Excludes the specialisations that the proof rules out.
\end{itemize}\medskip
\ifdefined\hypertarget\hypertarget{sourcepointer-1}{}\fi
\noindent\textit{Source reference.} Original lemma 02NW: focused-seeds/AGP/sources/stacks-snapshot-20261011/tags.txt, source lines 3495--3495. The lemma's own canonical tag, quoted so that readers can resolve the label inside the delivered text; it points at the counted statement itself and carries no proof content.\par\medskip
\ifdefined\hypertarget\hypertarget{sourcepointer-2}{}\fi
\noindent\textit{Source reference.} Original lemma 02ML: focused-seeds/AGP/sources/stacks-snapshot-20261011/tags.txt, source lines 3448--3448. Navigation pointer to a statement the counted proof refers to but does not contain; nothing load-bearing is delegated to it.\par\medskip
\ifdefined\hypertarget\hypertarget{sourcepointer-3}{}\fi
\noindent\textit{Source reference.} Original lemma 02NG: focused-seeds/AGP/sources/stacks-snapshot-20261011/tags.txt, source lines 3479--3479. Navigation pointer to a statement the counted proof refers to but does not contain; nothing load-bearing is delegated to it.\par\medskip
\ifdefined\hypertarget\hypertarget{sourcepointer-4}{}\fi
\noindent\textit{Source reference.} Original lemma 01ZV: focused-seeds/AGP/sources/stacks-snapshot-20261011/tags.txt, source lines 2610--2610. Navigation pointer to a statement the counted proof refers to but does not contain; nothing load-bearing is delegated to it.\par\medskip
\ifdefined\hypertarget\hypertarget{sourcepointer-5}{}\fi
\noindent\textit{Source reference.} Original lemma 02NE: focused-seeds/AGP/sources/stacks-snapshot-20261011/tags.txt, source lines 3477--3477. Navigation pointer to a statement the counted proof refers to but does not contain; nothing load-bearing is delegated to it.\par\medskip
\begin{lemma}
\label{lemma-generically-finite}
Let $X$, $Y$ be schemes.
Let $f : X \to Y$ be locally of finite type.
Let $\eta \in Y$ be a generic point of an irreducible component
of $Y$. The following are equivalent:
\begin{enumerate}
\item the set $f^{-1}(\{\eta\})$ is finite,
\item there exist affine opens $U_i \subset X$, $i = 1, \ldots, n$
and $V \subset Y$ with $f(U_i) \subset V$,
$\eta \in V$ and $f^{-1}(\{\eta\}) \subset \bigcup U_i$
such that each $f|_{U_i} : U_i \to V$ is finite.
\end{enumerate}
If $f$ is quasi-separated, then these are also equivalent to
\begin{enumerate}
\item[(3)] there exist affine opens $V \subset Y$,
and $U \subset X$ with $f(U) \subset V$,
$\eta \in V$ and $f^{-1}(\{\eta\}) \subset U$
such that $f|_U : U \to V$ is finite.
\end{enumerate}
If $f$ is quasi-compact and quasi-separated,
then these are also equivalent to
\begin{enumerate}
\item[(4)] there exists an affine open $V \subset Y$, $\eta \in V$
such that $f^{-1}(V) \to V$ is finite.
\end{enumerate}
\end{lemma}

\begin{proof}
The question is local on the base. Hence we may replace $Y$ by an
affine neighbourhood of $\eta$, and we may and do assume throughout
the proof below that $Y$ is affine, say $Y = \Spec(R)$.

\medskip\noindent
It is clear that (2) implies (1).
Assume that $f^{-1}(\{\eta\}) = \{\xi_1, \ldots, \xi_n\}$ is finite.
Choose affine opens $U_i \subset X$ with $\xi_i \in U_i$.
By Algebra, Lemma \ref{algebra-lemma-generically-finite} we see
that after replacing $Y$ by a standard open in
$Y$ each of the morphisms $U_i \to Y$ is finite.
In other words (2) holds.

\medskip\noindent
It is clear that (3) implies (1). Assume $f$ is quasi-separated and (1).
Write $f^{-1}(\{\eta\}) = \{\xi_1, \ldots, \xi_n\}$. There are no
specializations among the $\xi_i$ by Lemma \ref{lemma-finite-fibre}.
Since each $\xi_i$ maps to the generic
point $\eta$ of an irreducible component of $Y$, there cannot be a
nontrivial specialization $\xi \leadsto \xi_i$ in $X$
(since $\xi$ would map to $\eta$ as well).
We conclude each $\xi_i$ is a generic point of an
irreducible component of $X$.
Since $Y$ is affine and $f$ quasi-separated we see $X$ is quasi-separated.
By Properties, Lemma \ref{properties-lemma-maximal-points-affine}
we can find an affine open $U \subset X$ containing each $\xi_i$.
By Algebra, Lemma \ref{algebra-lemma-generically-finite} we see
that after replacing $Y$ by a standard open in
$Y$ the morphisms $U \to Y$ is finite.
In other words (3) holds.

\medskip\noindent
It is clear that (4) implies all of (1) -- (3) with no further assumptions
on $f$. Suppose that $f$ is quasi-compact and quasi-separated. We have to
show that the equivalent conditions (1) -- (3) imply (4).
Let $U$, $V$ be as in (3). Replace $Y$ by $V$. Since $f$ is quasi-compact
and $Y$ is quasi-compact (being affine) we see that $X$ is quasi-compact.
Hence $Z = X \setminus U$ is quasi-compact, hence the morphism
$f|_Z : Z \to Y$ is quasi-compact. By construction of $Z$ we see that
$\eta \not \in f(Z)$. Hence by
Lemma \ref{lemma-quasi-compact-generic-point-not-in-image}
we see that there exists an affine open
neighbourhood $V'$ of $\eta$ in $Y$ such that $f^{-1}(V') \cap Z = \emptyset$.
Then we have $f^{-1}(V') \subset U$ and this means
that $f^{-1}(V') \to V'$ is finite.
\end{proof}

\end{document}
